\begin{table}[t]
\centering
\small
\begin{adjustbox}{max width=\linewidth}
\begin{tabular}{rlrrrrrr}
\toprule
m & Variant & Max recall@10 & QPS @0.90 (recall, nprobe) & QPS @0.95 (recall, nprobe) & Bytes/vector in RAM & PQ part bytes/vector & PQ code bytes \\
\midrule
8 & plain IVF-PQ & 0.412 & not reached & not reached & 18.3 & 18.3 & 8 \\
8 & IVF-PQ + re-rank top 100 & 0.885 & not reached & not reached & 530.3 & 18.3 & 8 \\
16 & plain IVF-PQ & 0.590 & not reached & not reached & 26.3 & 26.3 & 16 \\
16 & IVF-PQ + re-rank top 100 & 0.985 & 4,508 (0.921, nprobe=32) & 2,917 (0.963, nprobe=64) & 538.3 & 26.3 & 16 \\
32 & plain IVF-PQ & 0.745 & not reached & not reached & 42.3 & 42.3 & 32 \\
32 & IVF-PQ + re-rank top 100 & 1.000 & 3,014 (0.927, nprobe=32) & 1,841 (0.973, nprobe=64) & 554.3 & 42.3 & 32 \\
64 & plain IVF-PQ & 0.885 & not reached & not reached & 74.3 & 74.3 & 64 \\
64 & IVF-PQ + re-rank top 100 & 1.000 & 1,875 (0.927, nprobe=32) & 1,067 (0.973, nprobe=64) & 586.3 & 74.3 & 64 \\
\bottomrule
\end{tabular}
\end{adjustbox}
\caption{PQ compression on SIFT (nlist 4096, 8-bit codes): plain IVF-PQ vs the same index with exact re-ranking of the top 100 candidates. Single-thread QPS on the test queries (median of 3 runs). Max recall: test recall of the config with the highest tune recall. Configs are chosen on the tune queries (recall $\geq$ target); test recall is reported as measured and may fall slightly below the target. 'Bytes/vector in RAM' for re-rank includes the raw float32 vectors; 'PQ part' is the IVF-PQ index alone.}
\label{tab:part3}
\end{table}
